\chapter{Hardware Overview}
\label{chap:fpgasteel-hw-overview}

\section{Scope of this chapter}
\label{sec:fpgasteel-hw-scope}

\prjname{FPGAsteel} started out as a full copy of the \prjname{FPGAlix} repository, subsequently modified independently. What actually carries over, though, is almost entirely the hardware: the \term{Platform Designer} system this project starts from is \prjname{FPGAlix}'s own \repopath{soc_system.qsys}. Everything \prjname{FPGAlix}'s hardware chapter documents about it carries over unchanged and is not repeated here: the rationale for \term{Avalon-ST}/\term{mSGDMA} streaming, the \term{OV7670} sensor and its adapter board, \term{PCLK} generation and its reset problem, the camera acquisition core's state machine, and the rationale for reaching \term{DDR3} through a dedicated \term{SDRAM} port rather than the \term{L3} interconnect. See Hardware Overview (PART~I -- \prjname{FPGAlix}, Chapter~\ref{chap:fpgalix-hw-overview}).

Very little carries over on the software side. Only one boot-chain submodule (|u-boot-socfpga|) is shared between the two repositories. \prjname{FPGAlix}'s own software, under its \repopath{sw/} folder (\cmd{camera_driver_2}, \cmd{hexip_daemon}, \cmd{vision_core}, embedded-\term{Linux} userspace code), has no counterpart here: \prjname{FPGAsteel}'s \repopath{sw/} folder is an unrelated, bare-metal example suite, documented starting from \nameref{chap:fpgasteel-board-bringup}.

The following sections document only what changed on top of that starting point: a handful of renamed instances on the camera capture path, a narrower capture datapath with \term{mSGDMA} driven by a hardware descriptor prefetcher instead of software-managed descriptors, and a full \term{VGA} output path with no counterpart in \prjname{FPGAlix}'s hardware (Chapter~\ref{chap:fpgalix-hw-overview}).

The hardware described here is common to every example under \repopath{sw/}: none of them requires recompiling it, regenerating its bitstream, or rebuilding the bootloader; only the software differs between them.

\section{Camera capture path}
\label{sec:fpgasteel-camera-capture}

\subsection{Renamed instances}
\label{subsec:fpgasteel-renamed-instances}

One of the first changes made with respect to \prjname{FPGAlix}'s \term{Platform Designer} system is renaming every \term{IP} instance on the camera capture path. Table~\ref{tab:fpgasteel-renamed-instances} reports the changes.

\begin{table}[htbp]
\centering
\begin{tabularx}{\textwidth}{@{} X X @{}}
\toprule
\textbf{FPGAlix instance} & \textbf{FPGAsteel instance} \\
\midrule
\texttt{\seqsplit{ov7670\_data\_interface\_0}} & \texttt{\seqsplit{ov7670\_video\_interface}} \\
\lightmidrule
\texttt{\seqsplit{pclk\_reset\_controller\_0}} & \texttt{\seqsplit{ov7670\_PCLK\_reset\_controller}} \\
\lightmidrule
\texttt{\seqsplit{mm\_clock\_crossing\_bridge\_0}} & \texttt{\seqsplit{ov7670\_AvalonMM\_clock\_bridge}} \\
\lightmidrule
\texttt{\seqsplit{dc\_fifo\_0}} & \texttt{\seqsplit{ov\_7670\_video\_dc\_fifo}} \\
\lightmidrule
\texttt{\seqsplit{st\_adapter\_0}} & \texttt{\seqsplit{ov7670\_st\_adapter}} \\
\lightmidrule
\texttt{\seqsplit{sc\_fifo\_0}} & Removed \\
\bottomrule
\end{tabularx}
\caption{Renamed and removed \term{Platform Designer} instances on the camera capture path.}
\label{tab:fpgasteel-renamed-instances}
\end{table}

All these instances remain the same type, except for the custom \term{IP} \hw{ov7670_data_interface}, for which a hardware update was made. Here, the frame size is no longer bound to the value of the \ccode{FRAME_SIZE} \term{VHDL} generic set during hardware compilation, but can now be changed by writing the \ccode{frame_size} register, accessible over the \term{Avalon-MM} bus. To preserve backward compatibility, the \ccode{FRAME_SIZE} generic remains settable, but it now only determines the value loaded into the \ccode{frame_size} register at reset. This is what lets the capture path switch between the \term{OV7670}'s 1-byte-per-pixel (\term{RAW}/processed \term{Bayer}) and 2-byte-per-pixel (\term{RGB565}/\term{RGB555}/\term{RGB444}/\term{YUV422}) output formats without recompiling the hardware.

\subsection{mSGDMA configuration}
\label{subsec:fpgasteel-msgdma-config}

Two further changes distinguish the capture path from \prjname{FPGAlix}'s:

\begin{itemize}
    \item \textbf{Width.} \prjname{FPGAlix}'s \hw{st_adapter_0} widens the camera's 8-bit \term{Avalon-ST} stream to 128 bit. This width reflects neither a bandwidth requirement nor the adapter's own maximum capability: with one of the \term{HPS}'s DRAM controller ports already occupied, 128 bit was the widest configuration available at the time, adopted because the fabric had sufficient resources to accommodate it.

    \prjname{FPGAsteel}'s \hw{ov7670_st_adapter} widens the same stream to 64 bit only, because adding the \term{VGA} output path to the same fabric leaves less room to spend on this width.

    The capture path itself would have worked with less than 64 bit, even 8 bit would have carried the required throughput, but going narrower would not have saved any resource: the dedicated \term{FPGA}-to-\term{HPS} \term{SDRAM} port each \term{mSGDMA} uses is provisioned the same way for any \term{Avalon-MM} interface up to 64 bit.\footnote{A 32- or 64-bit \term{Avalon-MM} interface, read-only or write-only, costs exactly one command port and one data port; only at 128 bit does the port cost double \cite[Table 12-3, p.~12-6]{CV-TRM}. The underlying data ports are natively 64-bit regardless of the configured interface width, so no narrower configuration costs less \cite[Table 12-2, p.~12-5]{CV-TRM}.} The \term{mSGDMA} \term{IP}'s 512-beat maximum burst count itself is unchanged; since the largest single burst into \term{DDR3} follows directly from it, it drops from 8 \term{KiB} at 128 bit (\prjname{FPGAlix}) to 4 \term{KiB} at 64 bit (\prjname{FPGAsteel}).
    \item \textbf{Descriptor management.} \prjname{FPGAsteel} enables the \term{mSGDMA}'s hardware descriptor prefetcher, where \prjname{FPGAlix} left it disabled and managed descriptors from software instead. The platform each project runs on drives this difference: \prjname{FPGAlix} runs embedded \term{Linux}, where using the prefetcher would have required developing a dedicated kernel driver for it. \prjname{FPGAsteel} is bare-metal, with direct access to the hardware, so enabling the prefetcher doesn't cost a \term{Linux} kernel integration, only direct register writes.

    \hw{ov7670_video_msgdma} fetches successive descriptors on its own out of a dedicated on-chip \term{RAM}, \hw{ov7670_prefercher_descriptor_mem}, re-arming itself between frames without \term{CPU} intervention.

    \prjname{FPGAlix}'s \hw{sc_fifo_0}, a buffer \term{FIFO} placed there to mitigate frame loss when memory and \term{CPU} were under load, has no counterpart here. This is not simply because descriptors now live on-chip rather than in \term{DDR3}; it is because the whole re-arm path is faster and more deterministic than in \prjname{FPGAlix}: bare-metal software re-arms each descriptor from its interrupt handler with lower, more predictable latency than a \term{Linux} driver would, and the on-chip descriptor memory itself is read directly, at \term{SRAM} speed, without \term{DDR3}'s burst-oriented access and its controller's \term{MPFE} arbitration overhead. Together, these remove the latency \hw{sc_fifo_0} existed to absorb. How the prefetcher itself operates (descriptor layout, ping-pong buffering) is documented in full in Paragraph~\ref{subsec:fpgasteel-descriptors-prefetcher} and is not repeated here.
\end{itemize}

\section{VGA output path (new)}
\label{sec:fpgasteel-vga-output}

\prjname{FPGAsteel} adds a second, symmetric datapath in the \term{FPGA} fabric, \term{Avalon-MM}$\rightarrow$\term{ST} this time, that reads frames back out of \term{DDR3} and drives the \term{DE1-SoC}'s \term{VGA} port directly, mirroring the capture path in reverse (software side: Paragraph~\ref{subsec:fpgasteel-simple-forwarding-vga-output}). None of the components below exist in \prjname{FPGAlix} \term{Platform Designer}'s system:

\begin{itemize}
    \item \textbf{\hw{vga_pll}} (kind |altera_pll|): generates a fixed 25 \term{MHz} pixel clock from \hw{clk_0} (the board's 50 \term{MHz} oscillator, the same \hw{sys_clk} source everything else in the design uses).
    \item \textbf{\hw{vga_reset_syncronizer}} (kind |pll_lock_reset_sync|, a new custom component, source \repopath{hw/my_vhdl/qsys_components/pll_lock_reset_sync.vhd}): releases the reset for the \hw{vga_clk} domain only once the system reset is released and \hw{vga_pll} reports locked.
    \item \textbf{\hw{vga_msgdma}} (kind |altera_msgdma|, \term{MM}$\rightarrow$\term{ST} direction): the read-back counterpart of \hw{ov7670_video_msgdma}, equally prefetcher-enabled with its own on-chip descriptor \term{RAM}, \hw{vga_prefercher_descriptor_mem}.
    \item \textbf{\hw{vga_dc_fifo}} (kind |altera_avalon_dc_fifo|): crosses the stream from \hw{sys_clk} into the \hw{vga_clk} domain immediately after \hw{vga_msgdma}'s \term{Avalon-ST} source port, while the stream is still kept at 64 bit wide.
    \item \textbf{\hw{vga_64_to_24_repacker}} (kind |avalon_st_64_to_24_repacker|, new custom \term{IP}, source \repopath{hw/my_vhdl/qsys_components/avalon_st_64_to_24_repacker.vhd}): named for the color depth it carries, \term{RGB888}, not for its output wire width. \hw{vga_controller}'s \term{Avalon-ST} sink is fixed by the Intel University Program \term{IP} at 30 bit (\ccode{dataBitsPerSymbol}=10, \ccode{symbolsPerBeat}=3, \term{RGB101010}), and \term{Platform Designer}'s automatic width adapter cannot bridge an 8-bit-symbol stream to a 10-bit-symbol one, since the symbol width itself changes and not just the beat count. This component narrows the stream from 64 bit down to 24 bit, then zero-extends each 8-bit \term{RGB888} channel into its own 10-bit slot (2 zero low-order bits per channel), producing exactly the 30-bit \term{RGB101010} beat \hw{vga_controller} expects, one pixel per beat, entirely inside the \hw{vga_clk} domain, after the crossing above.
    \item \textbf{\hw{vga_controller}} (kind |altera_up_avalon_video_vga_controller|, standard Altera University Program \term{IP}): drives the \term{DE1-SoC}'s \term{VGA} connector's sync and color signals at a fixed $640\times480$.
    \item A third dedicated \textbf{\term{FPGA}-to-\term{HPS} \term{SDRAM} port.} \prjname{FPGAlix} uses two: \hw{f2h_sdram0_data}, wired to a \term{JTAG} \term{Avalon} master for debug access, and \hw{f2h_sdram1_data} for its camera-capture \term{mSGDMA}. \prjname{FPGAsteel} simply adds a third, \hw{f2h_sdram2_data} port, for \hw{vga_msgdma}, for the same reason \prjname{FPGAlix} gives its own \term{mSGDMA} a dedicated port instead of routing through the \term{L3} interconnect (PART~I -- \prjname{FPGAlix}, Paragraph~\ref{sec:fpgalix-sdram-port}).
\end{itemize}

\section{System block diagram}
\label{sec:fpgasteel-system-block-diagram}

Figure~\ref{fig:fpgasteel-system-overview} shows \prjname{FPGAsteel}'s \repopath{soc_system.qsys}, redrawn from PART~I -- \prjname{FPGAlix}, Figure~\ref{fig:fpgalix-system-overview}. As described in Section~\ref{sec:fpgasteel-vga-output}, the \term{VGA} datapath, its hardware prefetcher, and its on-chip descriptor memory have been added; compare to \prjname{FPGAlix}, the camera capture path keeps its original layout, renamed per Paragraph~\ref{subsec:fpgasteel-renamed-instances} and re-labelled 64 bit per Paragraph~\ref{subsec:fpgasteel-msgdma-config}, with \hw{sc_fifo_0} removed and an on-chip descriptor \term{RAM} added for \hw{ov7670_video_msgdma}'s own prefetcher (Paragraph~\ref{subsec:fpgasteel-msgdma-config} explains why the former is no longer needed).

Clock generation is not drawn: \hw{vga_pll} and \hw{vga_reset_syncronizer} carry no data, they only stand up the \hw{vga_clk} domain the datapath runs in downstream of \hw{vga_msgdma}. That domain's color still appears in the figure, on \hw{vga_dc_fifo}, \hw{vga_64_to_24_repacker}, and \hw{vga_controller}, and the legend names \hw{vga_pll} as its source; both components are described in full in Section~\ref{sec:fpgasteel-vga-output}.

\begin{landscape}
\begin{figure}[p]
    \centering
    \includesvg[width=\linewidth,height=\textwidth,keepaspectratio]{FPGAsteelQSysDiagram}
    \caption{\texttt{soc\_system.qsys} block diagram, redrawn from PART~I -- \prjname{FPGAlix}, Figure~\ref{fig:fpgalix-system-overview} with the changes from Section~\ref{sec:fpgasteel-camera-capture} and Section~\ref{sec:fpgasteel-vga-output} applied.}
    \label{fig:fpgasteel-system-overview}
\end{figure}
\end{landscape}
